\chapter*{Введение}
\addcontentsline{toc}{chapter}{Введение}
\label{ch:intro}

\textbf{WordPress} — наиболее распространённая в мире система управления
контентом (CMS), используемая для публикации сайтов, блогов и
веб-приложений. Для её работы требуется стек \textbf{LAMP}: операционная
система Linux, веб-сервер Apache, СУБД MariaDB и интерпретатор PHP.

Цель работы --- развернуть и настроить WordPress на базе стека LAMP
на отдельной виртуальной машине в локальной сети, добавить
DNS-запись для нового сервера и убедиться в доступности сайта
с клиентской машины.

Задачи:
\begin{enumerate}
  \item Добавить DNS-запись \texttt{A} и \texttt{PTR} для ВМ
        \texttt{wordpress} на сервере BIND9 (\texttt{gateway}).
  \item Установить стек LAMP на ВМ \texttt{wordpress
        (192.168.\cfgLabStudentN.6)}.
  \item Создать базу данных и пользователя MariaDB для WordPress.
  \item Скачать, настроить и развернуть WordPress в
        \texttt{/var/www/html}.
  \item Настроить Virtual Host Apache и выполнить интерактивную
        установку WordPress через браузер.
  \item Проверить работу сайта и создать тестовую запись.
\end{enumerate}

\textbf{Стенд:}
\begin{itemize}
  \item ВМ \texttt{gateway} --- Ubuntu~22.04~LTS (VirtualBox),
        два адаптера: \texttt{enp0s3}~(WAN/NAT) и
        \texttt{enp0s8}~(LAN~Internal,~\texttt{\cfgGwIP/24});
        запущены BIND9 и isc-dhcp-server.
  \item ВМ \texttt{wordpress} --- Ubuntu~22.04~LTS~Server, один адаптер
        \texttt{enp0s3}~(Internal,~\texttt{\cfgWpIP/24});
        стек LAMP, WordPress.
  \item ВМ \texttt{desktop1} --- Ubuntu~22.04~Desktop, один адаптер
        \texttt{enp0s3}~(Internal), IP по DHCP;
        браузер Firefox для проверки.
\end{itemize}

\endinput
